\begin{figure}[htbp]
\centering
\begin{tikzpicture}[
  x=1cm,y=1cm,
  every node/.style={font=\small},
  box/.style={draw=linkblue!80!black,rounded corners=2pt,
    text width=5.9cm,align=center,inner sep=8pt,minimum height=11mm},
  branch/.style={box,text width=2.70cm,inner sep=6pt,minimum height=15mm},
  heading/.style={font=\small\bfseries,align=center},
  reduction/.style={-{Latex[length=2.1mm]},thick,draw=linkblue!80!black},
  implication/.style={-{Latex[length=2.1mm]},thick,dashed,draw=black!70}
]
\node[heading] at (0,0) {归约};
\node[heading] at (7.4,0) {NO 输入上的可靠性};

\node[box] (parity) at (0,-1.0)
  {$\mathrm{3SAT}$\\$\longrightarrow$ 带权奇偶间隙实例};
\node[box] (matrix) at (0,-2.8)
  {潜在矩阵测试\\
   稳定非线性映射：Lemma~\ref{lem:latent}};
\node[box] (ugc) at (0,-6.2)
  {无权简单二部\\唯一游戏\\
   Theorem~\ref{thm:ugc}};
\draw[reduction] (parity) -- (matrix);
\draw[reduction] (matrix) -- node[midway,fill=white,inner sep=4pt,
    text width=5.0cm,align=center]
  {取整与细分；\\并行重复} (ugc);

\node[branch] (assumption) at (5.65,-1.0)
  {假设测试接受率\\至少 $0.99$};
\node[branch] (no) at (9.15,-1.0)
  {NO 奇偶\\实例};
\node[branch] (decoding) at (5.65,-3.3)
  {解码后的建议策略\\一致概率至少 $\gamma>0$\\
   Proposition~\ref{prop:matrix-decoding}};
\node[branch] (clean) at (9.15,-3.3)
  {每一对策略：\\一致概率至多\\$\exp(-c k^{1/3})$\\
   Lemma~\ref{lem:clean}};
\node[box] (contradiction) at (7.4,-6.2)
  {对足够大的 $k$ 得到矛盾：\\
   $\exp(-c k^{1/3})<\gamma$};
\draw[implication] (assumption) -- (decoding);
\draw[implication] (no) -- (clean);
\draw[implication] (decoding) -- (contradiction);
\draw[implication] (clean) -- (contradiction);
\end{tikzpicture}
\caption{实线箭头描述归约；虚线箭头描述其可靠性论证。解码下界来自
高测试接受率，而干净坐标上界来自 NO 奇偶间隙，并适用于每个建议
策略。常数 $c,\gamma>0$ 在元组长度 $k$ 增长之前固定。}
\label{fig:overview}
\end{figure}
